\chapter{Quand l'échafaudage ne suffit plus : le contrat nu}

\begin{quote}
Vous n'avez pas besoin de ce chapitre pour écrire un jeu. Trois fonctions
suffisent, et c'était le chapitre~2. Celui-ci est pour le jour où votre état de
partie ne rentre plus dans la structure \texttt{Partie} --- une liste qui
grandit, un cache, un arbre. Ce jour-là, vous écrivez les vingt-quatre entrées
vous-même, et tout ce qui suit vous servira.
\end{quote}

Il se lira d'ailleurs beaucoup mieux maintenant que vous avez vu les dix-huit
fonctions à l'œuvre.

Ce chapitre construit un moteur complet, de la première ligne à la dernière. Le
fichier existe, il compile, et l'atelier le charge :
\nkfile{Applications/ConquerorLab/exemples/rules/RegleContratNu.cpp}.

Il joue \textbf{exactement le même jeu} que
\nkfile{exemples/rules/RegleFacile.cpp} du chapitre~2 : 570 lignes au lieu de
110, et pas une règle de plus. Comparer les deux est l'exercice le plus
instructif du cours.

Ouvrez-le à côté de ce texte. Nous le parcourons dans l'ordre, en expliquant à
chaque étape \emph{pourquoi} elle est là et \emph{ce qui casse} si on l'oublie.

\section{Ce que ce moteur fait}

\begin{nkcode}{RegleContratNu.cpp — resume}
plateau         carre 5x5, 4 voisins  (pas d'hexagone : une chose a la fois)
action          DUPLIQUER uniquement, portee 1
transformation  les ennemis voisins de la case OU L'ON VIENT DE POSER
                changent de camp -- et eux seuls
fin             des qu'un joueur ne peut plus dupliquer
decompte        le plus de totems gagne
\end{nkcode}

C'est déjà un jeu jouable, et c'est assez pour que l'atelier fasse \textbf{tout}
ce qu'il sait faire : partie humaine, IA, journal, rejeu, campagne de mesure.

\section{Les vingt-trois fonctions, par famille}

\texttt{NkcRulesVTable} compte vingt-trois entrées depuis l'ABI 3. Elles se
rangent en cinq familles, et cette carte suffit à ne jamais se perdre.

\begin{center}
\begin{tabular}{@{}p{2.9cm}p{6.4cm}p{4.5cm}@{}}
\toprule
\textbf{Famille} & \textbf{Fonctions} & \textbf{Rôle} \\
\midrule
cycle de vie & \texttt{Create}, \texttt{Destroy}, \texttt{CreateState},
  \texttt{DestroyState}, \texttt{CloneState}, \texttt{Setup} & fabriquer un
  moteur, puis des parties \\
réglages & \texttt{GetParamsSchemaJson}, \texttt{SetParam}, \texttt{GetParam},
  \texttt{LoadBoardJson}, \texttt{GetBoardJson} & ce qui change sans recompiler \\
jeu & \texttt{GetView}, \texttt{GenerateLegalMoves}, \texttt{IsLegalMove},
  \texttt{ApplyMove}, \texttt{IsFinished}, \texttt{GetWinner},
  \texttt{IsPlayerBlocked} & les règles proprement dites \\
sérialisation & \texttt{SerializeState}, \texttt{DeserializeState},
  \texttt{HashState} & rejeu, thread d'IA, diagnostic \\
géométrie \textsc{(abi 3)} & \texttt{GetCellCenter}, \texttt{GetCellShape} &
  \textbf{optionnel} — chapitre 5 \\
\bottomrule
\end{tabular}
\end{center}

\begin{nkretenir}
Convention de retour, partout : \textbf{0 = échec, 1 = succès}. Une fonction qui
échoue ne doit \textbf{rien} avoir modifié.
\end{nkretenir}

\section{La mémoire : jamais de \texttt{new} ni de \texttt{delete} bruts}

\begin{nkcode}{exemples/rules/RegleContratNu.cpp — section 1}
NkcAllocFn gAlloc = nullptr;
NkcFreeFn  gFree  = nullptr;

void *RawAlloc(usize n) { return gAlloc ? gAlloc(n) : std::malloc(n); }
void  RawFree(void *p)  { if (!p) return; if (gFree) gFree(p); else std::free(p); }
\end{nkcode}

L'atelier peut vous injecter ses propres allocateurs (NKMemory) par
\nkfile{nkc_rules_set_allocator}. Les deux chemins doivent exister : avec
injection, et sans.

\begin{nkpiege}[Ne jamais mélanger deux tas]
C'est une règle générale du dépôt : allouer avec l'allocateur du moteur et
libérer avec \texttt{std::free} (ou l'inverse) corrompt le tas — sous Windows,
cela donne un \texttt{c0000374} sans pile utile, souvent très loin du vrai
coupable. Le couple \texttt{RawAlloc}/\texttt{RawFree} garantit la symétrie.
\end{nkpiege}

\section{Les paramètres : aucune constante en dur}

\begin{nkcode}{exemples/rules/RegleContratNu.cpp — section 2}
enum ParamId : int32 {
    P_PORTEE = 0,     // distance maximale de duplication
    P_MAX_TOURS,      // garde-fou de simulation, jamais une regle de jeu
    P_COUNT
};

struct ParamDesc { const char *key, *group; int32 def, lo, hi; };

const ParamDesc kParams[P_COUNT] = {
    {"portee_duplication", "Duplication",   1,  1,   3},
    {"max_tours",          "Fin de partie", 200, 10, 100000},
};
\end{nkcode}

Une seule table décrit tout : la clé, le groupe d'affichage, la valeur par
défaut, les bornes. \textbf{Ajouter un réglage = ajouter une ligne.} Voici
pourquoi cela suffit à faire apparaître un champ à l'écran :

\begin{nkcode}{exemples/rules/RegleContratNu.cpp — V\_GetParamsSchemaJson}
k = std::snprintf(w, left,
                  "%s{\"key\":\"%s\",\"group\":\"%s\",\"type\":\"int\","
                  "\"min\":%d,\"max\":%d,\"def\":%d,\"val\":%d}",
                  i ? "," : "", kParams[i].key, kParams[i].group,
                  kParams[i].lo, kParams[i].hi, kParams[i].def, R->params[i]);
\end{nkcode}

Le panneau \emph{Règles} de l'atelier est \textbf{entièrement auto-généré} à
partir de cette chaîne. Il n'y a pas une ligne d'interface par paramètre.

\begin{center}
\begin{tabular}{@{}llp{5.6cm}@{}}
\toprule
\texttt{"type"} & \textbf{Widget} & \textbf{Champs supplémentaires} \\
\midrule
\texttt{"int"}  & champ numérique & \texttt{min}, \texttt{max} \\
\texttt{"bool"} & case à cocher & — \\
\texttt{"enum"} & liste déroulante & \texttt{"values":["A","B",\ldots]} \\
\bottomrule
\end{tabular}
\end{center}

Deux champs optionnels partout : \texttt{"label"} (sinon dérivé de la clé) et
\texttt{"help"} (affiché en infobulle).

\subsection{Borner, ne pas refuser}

\begin{nkcode}{exemples/rules/RegleContratNu.cpp — V\_SetParam}
int32 v = static_cast<int32>(value < 0 ? value - 0.5 : value + 0.5);
if (v < kParams[i].lo) v = kParams[i].lo;
if (v > kParams[i].hi) v = kParams[i].hi;
R->params[i] = v;
return 1;
\end{nkcode}

\begin{nkretenir}[Le banc d'essai vérifie ce comportement]
\begin{nkcode}{tests/NkcAbiHarness.cpp:121-124}
// Hors plage : doit etre borne au minimum du parametre (10), pas accepte.
R.SetParam(ri, "max_tours", 3);
CHECK(R.GetParam(ri, "max_tours") == 10.0,
      "une valeur hors plage est bornee, pas acceptee telle quelle");
\end{nkcode}
Le panneau \emph{Règles} relit le schéma \textbf{à chaque image} et n'a aucun
cache : si vous bornez, l'utilisateur voit tout de suite la valeur bornée. Un
cache local côté interface mentirait.
\end{nkretenir}

Notez que \texttt{SetParam} prend un \texttt{float64} — c'est le \textbf{seul}
flottant de la logique du contrat, et il ne la traverse pas : c'est un canal de
réglage, converti en entier dès la première ligne.

\section{L'état : votre représentation, leur vue}

\begin{nkcode}{exemples/rules/RegleContratNu.cpp — section 4}
struct State {
        NkcCellView cells[kCells];
        uint8       playerCount = 2;
        uint8       current     = 0;
        uint8       finished    = 0;
        int8        winner      = -2;
        uint32      turn        = 0;
        int32       energy[kMaxPlayers]         = {};
        int32       conquestTenths[kMaxPlayers] = {};
        int32       totemCount[kMaxPlayers]     = {};
        uint64      rng = 0x9E3779B97F4A7C15ull;   // PRNG PORTE PAR L'ETAT
};
\end{nkcode}

Le contrat n'impose rien de cette structure. Une seule contrainte, mais elle est
dure.

\begin{nkpiege}[L'état doit se copier tel quel]
\texttt{CloneState} est le \textbf{chemin chaud de l'IA} : appelé des milliers de
fois par seconde. Le contrat dit : « doit être rapide et ne jamais allouer ».
D'où :

\begin{nkcode}{RegleContratNu.cpp — V\_CloneState}
void V_CloneState(NkcRules, NkcState dst, const NkcState src) {
    if (dst && src) *static_cast<State *>(dst) = *static_cast<const State *>(src);
}
\end{nkcode}

Une seule affectation, parce que \texttt{State} est un POD. Si vous y mettez un
conteneur qui alloue, cette ligne alloue aussi, et votre IA perd un ordre de
grandeur.
\end{nkpiege}

La même contrainte explique \texttt{SerializeState} :

\begin{nkcode}{RegleContratNu.cpp — V\_SerializeState}
uint32 V_SerializeState(NkcRules, const NkcState st, void *buf, uint32 cap) {
    const uint32 need = static_cast<uint32>(sizeof(State));
    if (!buf || cap < need) return need;   // renvoie la taille NECESSAIRE
    std::memcpy(buf, st, need);
    return need;
}
\end{nkcode}

\begin{nkpiege}[Le protocole en deux temps]
Appelée avec \texttt{buf == nullptr}, la fonction doit renvoyer la \textbf{taille
nécessaire} sans rien écrire. C'est ainsi que l'atelier dimensionne son tampon
avant de transférer l'état vers le thread d'IA. Renvoyer 0 dans ce cas fait
échouer la réflexion — et l'atelier affichera « SerializeState annonce une taille
nulle » dans la barre du plateau.
\end{nkpiege}

\section{Générer les coups}

C'est la fonction la plus importante du fichier. Trois exigences s'y croisent.

\begin{nkcode}{exemples/rules/RegleContratNu.cpp — GenMoves}
for (int32 i = 0; i < kCells; ++i) {
    if (s->cells[i].owner != static_cast<int8>(me)) continue;
    const NkcCoord src = R->coords[i];

    for (int32 k = 0; k < nbr; ++k) {
        const NkcCoord dst = Neighbor(NkcTopology::Square4, src, k);
        const int32    di  = IndexOf(dst);
        if (di < 0) continue;
        if (s->cells[di].owner != kCellEmpty) continue;
        if (Distance(NkcTopology::Square4, src, dst) > range) continue;

        if (n < cap) {
            NkcMove m;
            std::memset(&m, 0, sizeof(m));
            m.kind        = NkcMoveKind::Duplicate;
            m.player      = me;
            m.from        = src;
            m.to          = dst;
            m.targetLevel = -1;
            m.powerId     = -1;
            out[n]        = m;
        }
        ++n;
    }
}
\end{nkcode}

\textbf{Exigence 1 — l'ordre est normatif.} Cases par index croissant, puis
voisins dans l'ordre de \texttt{Neighbor}. Deux exécutions produisent la même
liste, sur toute plateforme.

\textbf{Exigence 2 — la mise à zéro est obligatoire.}

\begin{nkpiege}[\texttt{memset} avant de remplir, toujours]
\begin{quote}\small
Coups. POD comparable octet à octet APRÈS mise à zéro : le module DOIT remplir
intégralement les champs inutilisés avec zéro.
\hfill\textsc{ConquerorRulesABI.h:112}
\end{quote}

\texttt{NkcMove} contient \texttt{NkcCoord fuseCells[12]}, inutilisés au palier 0.
Sans \texttt{memset}, ils contiennent ce que la pile y a laissé. Or
\texttt{IsLegalMove} compare les coups \textbf{octet à octet} : deux coups
identiques deviendraient différents, et l'atelier rejetterait le coup de votre
propre IA. Le symptôme est spectaculaire — « l'IA a proposé un coup ILLÉGAL » —
et la cause est trois lignes plus haut.
\end{nkpiege}

Ce même \texttt{fuseCells} est ce qui rend une fusion \textbf{jouable ou non à
la souris} : un coup \texttt{Fuse} laisse \texttt{from} à zéro et met ses cases
là. Ce que le joueur voit alors à l'écran, et où il doit cliquer, sont décrits au
chapitre \textbf{«~Jouer un coup à la souris~»}.

\textbf{Exigence 3 — \texttt{cap} peut être trop petit.} La fonction écrit
\emph{au plus} \texttt{cap} coups mais renvoie le nombre \emph{total}. Remarquez
que \texttt{++n} est en dehors du \texttt{if (n < cap)} : l'appelant apprend ainsi
qu'il doit agrandir son tampon.

\subsection{Le cas de PASSER}

\begin{nkcode}{exemples/rules/RegleContratNu.cpp — GenMoves, fin}
// PASSER n'est legal que si RIEN d'autre ne l'est (REGLES 6.2).
if (n == 0) {
    if (cap > 0) { /* ... un coup Pass ... */ }
    return 1;
}
\end{nkcode}

Ce n'est pas un détail de confort : c'est ce qui rend vraie l'équivalence
suivante, que le banc d'essai vérifie automatiquement.

\section{L'équivalence qui teste tout le reste}

\begin{nkcode}{tests/NkcAbiHarness.cpp:107-111}
const uint32 n = R.GenerateLegalMoves(ri, probe, buf, 512);
const bool onlyPass = (n == 1 && buf[0].kind == NkcMoveKind::Pass);
const bool blocked  = R.IsPlayerBlocked(ri, probe, pv.current) != 0;
if (onlyPass != blocked) { equiv = false; break; }
\end{nkcode}

Autrement dit : \textbf{« aucun coup légal » doit être exactement équivalent à
« joueur bloqué »}. Les règles le disent en toutes lettres :

\begin{quote}\small
Aucune action légale mais joueur non bloqué : impossible par construction. Le
moteur doit \textbf{asserter cette équivalence en test} : c'est le meilleur test
d'intégrité du générateur de coups.
\hfill\textsc{regles\_completes\_v2.md:467}
\end{quote}

\begin{nkretenir}[Pourquoi ce test attrape presque tout]
\texttt{GenerateLegalMoves} et \texttt{IsPlayerBlocked} sont deux chemins de code
indépendants qui répondent à la même question. S'ils divergent, l'un des deux a
tort — et c'est presque toujours le générateur, sur un cas de bord : case hors
plateau, case bloquée, portée mal comparée.

Dans \nkfile{RegleContratNu.cpp} les deux partagent délibérément la même
primitive, \texttt{CanPlay}, ce qui rend l'équivalence vraie par construction.
\end{nkretenir}

\section{Appliquer un coup, et la cascade}

\begin{nkcode}{exemples/rules/RegleContratNu.cpp — DoApply, placement}
// --- placement : la source RESTE INTACTE (REGLES 7.2) ---
s->cells[di].owner = static_cast<int8>(mv->player);
s->cells[di].level = 0;
Emit(sink, user, NkcEventKind::TotemDuplicated, mv->player, mv->from, mv->to, 0);
s->conquestTenths[mv->player] += 10;   // +1,0 PC, en DIXIEMES entiers
\end{nkcode}

Puis la transformation, et \textbf{c'est ici qu'on se trompe} :

\begin{nkcode}{exemples/rules/RegleContratNu.cpp — DoApply, transformation}
const int32 nbr     = NeighborCount(NkcTopology::Square4);
int32       flipped = 0;
for (int32 k = 0; k < nbr; ++k) {
    const int32 ni = IndexOf(Neighbor(NkcTopology::Square4, mv->to, k));
    if (ni < 0) continue;
    const int8 o = s->cells[ni].owner;
    if (o < 0 || o == static_cast<int8>(mv->player)) continue;
    /* ... retourner ce totem ... */
    ++flipped;
}
if (flipped >= 2)
    Emit(sink, user, NkcEventKind::Cascade, mv->player, mv->to, mv->to, flipped);
\end{nkcode}

\begin{nkpiege}[Une transformation ne déclenche pas de transformation]
\begin{quote}\small
Une transformation ne déclenche pas de nouvelle transformation. Seul le totem
\emph{placé} est un déclencheur. Les cascades naissent de la multiplicité des
voisins, pas d'une propagation. C'est le point le plus facile à implémenter de
travers.
\hfill\textsc{regles\_completes\_v2.md:246}
\end{quote}

La boucle ci-dessus parcourt \textbf{une seule fois} les voisins de
\texttt{mv->to}. Si vous la rappelez sur chaque case retournée, tout le plateau
bascule au premier coup et le jeu n'existe plus. La « cascade » que l'atelier
affiche en gros au centre de l'écran, c'est \texttt{flipped >= 2} — plusieurs
voisins d'un coup, pas une réaction en chaîne.
\end{nkpiege}

\subsection{Les événements : décrire, pas animer}

\texttt{Emit} pousse un \texttt{NkcEvent} vers un collecteur fourni par
l'appelant. Le moteur ne sait rien de l'animation : il décrit ce qui s'est
produit, la présentation décide comment le montrer.

Ce collecteur peut être \texttt{nullptr} — et il l'est presque toujours, parce
qu'une IA ne veut pas payer le coût des événements pendant ses simulations.
\textbf{Testez-le} : \texttt{Emit} commence par \texttt{if (!sink) return;}.

\begin{nkretenir}[Les événements servent aussi à l'aperçu]
Quand vous survolez une destination dans l'atelier, les totems qui seraient
retournés s'entourent de rouge. L'atelier ne le \emph{déduit} pas des règles : il
clone l'état, joue le coup pour de faux avec un collecteur, et lit les
\texttt{TotemTransformed}.

Conséquence : \textbf{votre aperçu reste juste même quand vous changez la règle de
transformation}, sans qu'on touche à l'interface.
\end{nkretenir}

\section{Fin de partie et garde-fou}

\begin{nkcode}{exemples/rules/RegleContratNu.cpp — CheckEnd}
// Garde-fou de SIMULATION : sans lui, une partie sur dix mille ne finit
// jamais et le batch tourne pour toujours (REGLES 12.3).
if (!over && static_cast<int32>(s->turn) >=
                 R->params[P_MAX_TOURS] * static_cast<int32>(s->playerCount))
    over = true;
\end{nkcode}

\begin{nkpiege}[\texttt{max\_tours} n'est PAS une règle de jeu]
\begin{quote}\small
\texttt{max\_tours} n'est pas une règle de jeu : c'est une sécurité pour que
10\,000 parties IA-vs-IA se terminent toujours. Le rapport de batch doit indiquer
combien de parties ont fini par ce garde-fou. \textbf{Un taux non nul signale une
pathologie des règles, pas un réglage à monter.}
\hfill\textsc{regles\_completes\_v2.md:451}
\end{quote}

Le panneau \emph{Métriques} affiche ce taux dans une tuile à part, rouge dès
qu'il dépasse zéro. Si vous le voyez monter, ne montez pas \texttt{max\_tours} :
cherchez pourquoi vos parties ne se terminent pas.
\end{nkpiege}

\section{Les deux symboles exportés}

\begin{nkcode}{exemples/rules/RegleContratNu.cpp — fin de fichier}
NKC_MODULE_EXPORT void nkc_rules_set_allocator(NkcAllocFn a, NkcFreeFn f)
{ gAlloc = a; gFree = f; }

NKC_MODULE_EXPORT void nkc_rules_get_factory(NkcRulesFactory *out)
{ FillFactory(out); }
\end{nkcode}

\begin{nkpiege}[L'erreur du premier jour]
Les oublier donne, dans le panneau \emph{Modules} :

\begin{nkcode}{src/ConquerorLab/NkcModuleHost.h — BuildAndLoad}
e.log  = "Symbole introuvable : ";
e.log += symFactory;
e.log += " - as-tu bien termine ton fichier par la macro d'export ?";
\end{nkcode}

Le message a été écrit pour cette erreur précise, parce qu'elle arrive à tout le
monde une fois.
\end{nkpiege}

Et la version d'ABI, dans \texttt{FillFactory} :
\texttt{out->info.abiVersion = kRulesAbiVersion;}. L'atelier \textbf{refuse} un
module dont l'\texttt{abiVersion} diffère, avec un message qui donne les deux
numéros. Si vous lisez « ABI 2, attendue 3 », vous avez compilé contre de vieux
en-têtes.

\section{Votre premier module, en quatre gestes}

\begin{enumerate}
  \item copiez \nkfile{exemples/rules/RegleContratNu.cpp} vers
        \nkfile{travail/rules/mes\_regles.cpp} ;
  \item changez le nom dans \texttt{FillFactory} — sinon vous aurez deux entrées
        identiques dans le menu et vous ne saurez pas laquelle vous testez ;
  \item sauvegardez, attendez une seconde : le panneau \emph{Modules} affiche
        votre module ;
  \item sélectionnez-le, puis \textbf{Nouvelle partie}.
\end{enumerate}

À partir de là, modifiez \textbf{une chose à la fois} et mesurez. C'est tout le
métier.

\begin{nkexo}[1 — La portée]
\texttt{portee\_duplication} est réglable de 1 à 3, mais \texttt{GenMoves}
n'énumère que les \emph{voisins immédiats} : passer la portée à 2 ne change donc
rien. Corrigez-le. Indice : \texttt{Neighbor} ne donne que la distance 1 ; il vous
faut parcourir les cases et filtrer par \texttt{Distance(...) <= range}.
\end{nkexo}

\begin{nkexo}[2 — Casser l'équivalence exprès]
Dans \texttt{V\_IsPlayerBlocked}, remplacez \texttt{CanPlay(R, s, player)} par
\texttt{true}. Recompilez, lancez une partie IA contre IA. Que se passe-t-il ?
Reproduisez ensuite le test du banc d'essai à la main. Vous venez d'écrire votre
premier test d'intégrité.
\end{nkexo}

\begin{nkexo}[3 — Le \texttt{memset} manquant]
Retirez le \texttt{std::memset(\&m, 0, sizeof(m))} de \texttt{GenMoves},
recompilez, et faites jouer une IA. Décrivez ce que vous observez et
\textbf{expliquez-le} par le fonctionnement de \texttt{IsLegalMove}. Remettez
ensuite la ligne.
\end{nkexo}

\begin{nkexo}[4 — Le troisième joueur]
\texttt{V\_Setup} force \texttt{playerCount} à 2. Faites-en un moteur à trois
joueurs : placez un troisième totem de départ, faites tourner l'ordre de jeu, et
vérifiez le décompte en cas d'égalité. Que doit-il se passer quand un joueur n'a
plus aucun totem ?
\end{nkexo}
